% GENERATED FILE. Edit canonical lessons, then run npm run generate:books.
% canonical-source-hash: fnv1a64:34f4fa179c2cfdcb
% canonical-lessons: ZH-C215-songhuo, ZH-C215-qianshou, ZH-C215-jicun, ZH-C215-xiyidian, ZH-C215-ganxi

\chapter{Deliveries and the Laundry}
\label{ch:zh-deliveries-and-the-laundry}
\hlchaptermodality{215}

\begin{chapteropening}
\textbf{By the end of this chapter:} \emph{I can have goods delivered, sign for a parcel, leave luggage in safekeeping and use the laundry.}

Complete the last lesson of chapter 215: I can have goods delivered, sign for a parcel, leave luggage in safekeeping and use the laundry.
\end{chapteropening}

\section[sònghuò]{sònghuò — to deliver goods}

\label{lesson:ZH-C215-songhuo}



\begin{quote}
Before the new one: say the Mandarin for a way to deal with it, then the Mandarin for what should I do.
\end{quote}

\subsection*{You'll want to know: sònghuò}
\emph{sònghuò} — \textquotedblleft{}to deliver goods\textquotedblright{}.

\begin{quote}
\emph{tāmen míngtiān sònghuò} — They're delivering tomorrow.
\end{quote}

Its characters come later; for now it is a word for the ear.

\begin{cousinweb}
\emph{sònghuò} is built from \emph{sòng} 'to send' and \emph{huò} 'goods', both inherited from Middle Chinese, \emph{suwngH} and \emph{xwaH} (Guangyun, in Baxter's transcription), which glosses the first 'to send off'.
\end{cousinweb}

\begin{grammarlens}[title={What you've built}]
One more word for everyday actions.
\end{grammarlens}

\subsection*{Your turn}
\begin{itemize}
\raggedright
  \item \emph{Say it:} \emph{sònghuò}
  \item \emph{Say it:} \emph{sònghuò}, once more
  \item \emph{Say it:} \emph{tāmen míngtiān sònghuò} — They're delivering tomorrow.
  \item \emph{Recall:} say the Mandarin for a way to deal with it, then the Mandarin for what should I do, then say \emph{sònghuò} again
\end{itemize}

\begin{tcolorbox}[breakable,colback=teal!4,colframe=teal!35!black,title={Before you move on}]
What does sònghuò mean? (\textquotedblleft{}To deliver goods\textquotedblright{}.) Say it once more.
\end{tcolorbox}
\section[qiānshōu]{qiānshōu — to sign for (a delivery)}

\label{lesson:ZH-C215-qianshou}



\begin{quote}
Before the new one: say the Mandarin for what should I do, then the Mandarin for to deliver goods.
\end{quote}

\subsection*{You'll want to know: qiānshōu}
\emph{qiānshōu} — \textquotedblleft{}to sign for (a delivery)\textquotedblright{}.

\begin{quote}
\emph{kuàidì dào le qǐng qiānshōu} — Your delivery is here; please sign for it.
\end{quote}

Its characters come later; for now it is a word for the ear.

\begin{cousinweb}
\emph{qiānshōu} is built from \emph{qiān} 'to sign' and \emph{shōu} 'to receive'. The second is inherited from Middle Chinese \emph{syuw} (Guangyun, in Baxter's transcription); no dictionary we use records an older form of the first, so its history is not settled.
\end{cousinweb}

\begin{grammarlens}[title={What you've built}]
One more word for everyday actions.
\end{grammarlens}

\subsection*{Your turn}
\begin{itemize}
\raggedright
  \item \emph{Say it:} \emph{qiānshōu}
  \item \emph{Say it:} \emph{qiānshōu}, once more
  \item \emph{Say it:} \emph{kuàidì dào le qǐng qiānshōu} — Your delivery is here; please sign for it.
  \item \emph{Recall:} say the Mandarin for what should I do, then the Mandarin for to deliver goods, then say \emph{qiānshōu} again
\end{itemize}

\begin{tcolorbox}[breakable,colback=teal!4,colframe=teal!35!black,title={Before you move on}]
What does qiānshōu mean? (\textquotedblleft{}To sign for (a delivery)\textquotedblright{}.) Say it once more.
\end{tcolorbox}
\section[jìcún]{jìcún — to leave in safekeeping}

\label{lesson:ZH-C215-jicun}



\begin{quote}
Before the new one: say the Mandarin for to deliver goods, then the Mandarin for to sign for.
\end{quote}

\subsection*{You'll want to know: jìcún}
\emph{jìcún} — \textquotedblleft{}to leave in safekeeping\textquotedblright{}.

\begin{quote}
\emph{xíngli kěyǐ jìcún ma} — Can I leave my luggage here?
\end{quote}

Its characters come later; for now it is a word for the ear.

\begin{cousinweb}
\emph{jìcún} is built from \emph{jì} 'to entrust' and \emph{cún} 'to keep', both inherited from Middle Chinese, \emph{kjeH} and \emph{dzwon} (Guangyun, in Baxter's transcription), which glosses the first 'to entrust'.
\end{cousinweb}

\begin{grammarlens}[title={What you've built}]
One more word for everyday actions.
\end{grammarlens}

\subsection*{Your turn}
\begin{itemize}
\raggedright
  \item \emph{Say it:} \emph{jìcún}
  \item \emph{Say it:} \emph{jìcún}, once more
  \item \emph{Say it:} \emph{xíngli kěyǐ jìcún ma} — Can I leave my luggage here?
  \item \emph{Recall:} say the Mandarin for to deliver goods, then the Mandarin for to sign for, then say \emph{jìcún} again
\end{itemize}

\begin{tcolorbox}[breakable,colback=teal!4,colframe=teal!35!black,title={Before you move on}]
What does jìcún mean? (\textquotedblleft{}To leave in safekeeping\textquotedblright{}.) Say it once more.
\end{tcolorbox}
\section[\texorpdfstring{\zh{洗衣店} (xǐyīdiàn)}{洗衣店 (xǐyīdiàn)}]{\zh{洗衣店} (xǐyīdiàn) — laundry shop}

\label{lesson:ZH-C215-xiyidian}



\begin{quote}
Before the new one: say the Mandarin for to sign for, then the Mandarin for to leave in safekeeping.
\end{quote}

\subsection*{You'll want to know: \zh{洗衣店}}
\textbf{\zh{洗衣店}} — \emph{xǐyīdiàn} — \textquotedblleft{}laundry shop\textquotedblright{}.

\begin{quote}
\textbf{\zh{洗衣店在书店东边}} \emph{xǐyīdiàn zài shūdiàn dōngbian} — The laundry is to the east of the bookshop.
\end{quote}

\begin{cousinweb}
\textbf{\zh{洗衣店}} (\emph{xǐyīdiàn}) is built from \zh{洗} 'to wash', \zh{衣} 'clothes' and \zh{店} 'a shop', all inherited from Middle Chinese, \emph{sejX}, \emph{'j+j} and \emph{temH} (Guangyun, in Baxter's transcription).
\end{cousinweb}

\begin{grammarlens}[title={What you've built}]
One more word for everyday actions.
\end{grammarlens}

\subsection*{Your turn}
\begin{itemize}
\raggedright
  \item \emph{Say it:} \emph{xǐyīdiàn}
  \item \emph{Say it:} \emph{xǐyīdiàn}, once more
  \item \emph{Say it:} \emph{xǐyīdiàn zài shūdiàn dōngbian} — The laundry is to the east of the bookshop.
  \item \emph{Recall:} say the Mandarin for to sign for, then the Mandarin for to leave in safekeeping, then say \emph{xǐyīdiàn} again
\end{itemize}

\begin{tcolorbox}[breakable,colback=teal!4,colframe=teal!35!black,title={Before you move on}]
What does \zh{洗衣店} mean? (\textquotedblleft{}Laundry shop\textquotedblright{}.) Say it once more.
\end{tcolorbox}
\section[gānxǐ]{gānxǐ — to dry-clean}

\label{lesson:ZH-C215-ganxi}



\begin{quote}
Before the new one: say the Mandarin for to leave in safekeeping, then the Mandarin for laundry shop.
\end{quote}

\subsection*{You'll want to know: gānxǐ}
\emph{gānxǐ} — \textquotedblleft{}to dry-clean\textquotedblright{}.

\begin{quote}
\emph{zhège kěyǐ gānxǐ ma} — Can this be dry-cleaned?
\end{quote}

Its characters come later; for now it is a word for the ear.

\begin{cousinweb}
\emph{gānxǐ} is built from \emph{gān} 'dry' and \zh{洗} 'to wash', both inherited from Middle Chinese, \emph{kan} and \emph{sejX} (Guangyun, in Baxter's transcription).
\end{cousinweb}

\begin{grammarlens}[title={What you've built}]
One more word for everyday actions.
\end{grammarlens}

\subsection*{Your turn}
\begin{itemize}
\raggedright
  \item \emph{Say it:} \emph{gānxǐ}
  \item \emph{Say it:} \emph{gānxǐ}, once more
  \item \emph{Say it:} \emph{zhège kěyǐ gānxǐ ma} — Can this be dry-cleaned?
  \item \emph{Recall:} say the Mandarin for to leave in safekeeping, then the Mandarin for laundry shop, then say \emph{gānxǐ} again
\end{itemize}

\begin{tcolorbox}[breakable,colback=teal!4,colframe=teal!35!black,title={Before you move on}]
What does gānxǐ mean? (\textquotedblleft{}To dry-clean\textquotedblright{}.) Say it once more.
\end{tcolorbox}
